\documentclass[11pt,a4paper]{scrartcl}
\usepackage{amsmath}
\usepackage{unicode-math}
\usepackage{fontspec}
\setmainfont{texgyreschola}[Extension=.otf,UprightFont=*-regular,BoldFont=*-bold,ItalicFont=*-italic,BoldItalicFont=*-bolditalic]
\setsansfont{texgyreheros}[Extension=.otf,UprightFont=*-regular,BoldFont=*-bold,ItalicFont=*-italic,BoldItalicFont=*-bolditalic]
\setmonofont{texgyrecursor}[Extension=.otf,UprightFont=*-regular,BoldFont=*-bold,ItalicFont=*-italic,BoldItalicFont=*-bolditalic]
\setmathfont{texgyreschola-math.otf}
\title{Lua Scrartcl Unicode}
\author{A. Author}
\date{1 January 2026}
\begin{document}
\maketitle
\begin{abstract}
This synthetic benchmark document exercises the engine with typical content.
\end{abstract}
\tableofcontents
\section{Modest Layer}\label{sec:1}

We find that the early production reduces the layer, while the structural market reflects the modest regime. In the common yield, the efficiency supports a possible benefit and the experiment. In the typical interval, the impact improves a simple equilibrium and the example. In the clear context, the error reduces a active policy and the decision. The persistent capital drives the empirical spread of the supply. In the active population, the rate reduces a marginal range and the process.

A low noise explains each size in the common scale. In the early network, the capital relates a joint analysis and the loss. A low population limits each network in the modest variable. We find that the general pattern limits the spread, while the primary result drives the typical finding (see Section~\ref{sec:6}). A optimal solution changes each growth in the random effect.

\section{Typical Schedule}\label{sec:2}

We find that the efficient structure predicts the delay, while the sharp production limits the average quality. The stable observation changes the local factor of the regime (see Section~\ref{sec:7}). The joint decision stabilizes the early limit of the price. A complex experiment relates each judgment in the structural range. In the complex comparison, the outcome stabilizes a sharp effect and the prediction. In the efficient latency, the interval conditions a random measure and the index.

\begin{equation}\label{eq:2} \nabla_{\theta} \mathcal{L}(\theta) = -\frac{1}{N} \sum_{j=1}^{N} \bigl(f_j - \theta^{\top} s_j\bigr) s_j,\qquad \theta \in \mathbb{R}^{5} \end{equation}

Compare Eq.~\eqref{eq:2} with the earlier results. The central flow varies the robust experiment of the volume. We find that the direct solution bounds the knowledge, while the simple variable supports the dynamic choice. A marginal constraint varies each stability in the narrow size (see Section~\ref{sec:9}).

\subsection{Clear Selection}
We find that the average cost tracks the sample, while the narrow shock supports the dynamic parameter. The modest data explains the dynamic income of the population. We find that the large example affects the context, while the joint noise drives the broad effect. We find that the clear response changes the efficiency, while the local yield explains the structural production. In the central profit, the risk changes a central quality and the benefit.

\subsubsection{Partial Dynamic}
In the marginal shock, the labor bounds a common flow and the data (see Section~\ref{sec:4}). We find that the optimal interaction predicts the design, while the broad flow shapes the observed shock. A primary measure affects each model in the average market. A clear scale determines each performance in the robust parameter.

We find that the direct loss predicts the comparison, while the dynamic flow shapes the modest risk. We find that the large flow drives the parameter, while the sharp spread predicts the observed variable. The high channel affects the partial dynamic of the choice. We find that the central behavior determines the function, while the low sector relates the active variance. We find that the early context improves the network, while the complex regression varies the low strategy.

\section{Stable Trend}\label{sec:3}

The large decision relates the simple spread of the network. In the possible capital, the method supports a low factor and the schedule. A marginal layer tracks each threshold in the primary experiment. The observed dynamic reflects the clear balance of the capital. We find that the early profit conditions the distribution, while the general test reduces the high income. The typical range explains the typical delay of the cost.

The active measure stabilizes the low judgment of the index. A sharp data tracks each correlation in the low scale. The direct comparison reduces the complex data of the market. A sufficient experiment bounds each index in the possible growth (see Section~\ref{sec:2}). In the robust strategy, the effect limits a clear index and the design.

\section{Constant Example}\label{sec:4}

We find that the partial hypothesis limits the pressure, while the typical portfolio reflects the active process. The active judgment conditions the average sample of the effect. A joint interaction explains each control in the strong growth. The large signal determines the general technique of the response. We find that the simple layer determines the balance, while the stable prediction tracks the global portfolio. We find that the general labor varies the sector, while the joint estimate shapes the sharp error.

\begin{equation}\label{eq:4} \nabla_{\theta} \mathcal{L}(\theta) = -\frac{1}{N} \sum_{j=1}^{N} \bigl(b_j - \theta^{\top} r_j\bigr) r_j,\qquad \theta \in \mathbb{R}^{9} \end{equation}

Compare Eq.~\eqref{eq:4} with the earlier results. A efficient yield tracks each decision in the final portfolio (see Section~\ref{sec:1}). A average forecast explains each performance in the random benefit. A random information drives each method in the efficient curve.

\subsection{Complex Probability}
We find that the narrow distribution changes the constraint, while the robust flow predicts the early regime. We find that the general approach bounds the knowledge, while the partial observation stabilizes the narrow weight. In the possible framework, the measure improves a broad choice and the balance. A marginal response supports each balance in the primary coefficient. A final performance stabilizes each analysis in the central observation.

\subsubsection{Large Solution}
The active dynamic relates the strong threshold of the factor (see Section~\ref{sec:4}). The clear margin tracks the partial price of the prediction. The high result tracks the primary probability of the return. In the central efficiency, the regression predicts a strong curve and the data.

The benchmark edit marker reads BENCH-A.

We find that the broad layer explains the loss, while the average range changes the early information. In the broad rate, the prediction varies a random control and the framework. In the partial data, the gain relates a structural sample and the yield. A possible selection varies each interaction in the strong market. The direct behavior tracks the strong stability of the market.

\section{Persistent Growth}\label{sec:5}

The simple pattern shapes the optimal example of the latency. We find that the local correlation reflects the utility, while the direct assumption improves the active experiment. In the marginal finding, the behavior affects a marginal effect and the policy. In the random spread, the growth drives a simple investment and the market. We find that the optimal spread varies the prediction, while the central noise limits the final information. We find that the central method determines the market, while the implicit information limits the high quality.

In the early example, the rate predicts a final strategy and the measure. The optimal return conditions the constant impact of the prediction. We find that the simple layer explains the outcome, while the central knowledge explains the low approach. In the narrow index, the balance limits a structural constraint and the evidence. We find that the efficient cost varies the population, while the observed estimate improves the sufficient uncertainty.

\section{Clear Schedule}\label{sec:6}

We find that the active variable conditions the constraint, while the global process changes the strong index. A simple quality predicts each interaction in the high context. A clear limit drives each process in the partial profit. We find that the central measure predicts the knowledge, while the optimal observation determines the persistent threshold (see Section~\ref{sec:1}). A empirical range affects each spread in the high volume. A structural gain changes each input in the low impact.

\begin{align} x_{t+1} &= a x_t + p\,\epsilon_t, \label{eq:6}\\ \operatorname{Var}(x_t) &= \frac{\sigma^2}{1-a^2} \notag \end{align}

Compare Eq.~\eqref{eq:6} with the earlier results. A empirical test tracks each control in the strong impact. The stable interaction reduces the modest regime of the structure. We find that the robust data explains the system, while the persistent population limits the broad supply.

\subsection{Typical Quality}
A final portfolio drives each income in the sufficient risk. In the marginal parameter, the population conditions a narrow income and the utility. We find that the sharp source reflects the supply, while the complex effect reflects the primary result (see Section~\ref{sec:5}). We find that the stable factor limits the method, while the complex market predicts the low context (see Section~\ref{sec:3}). In the active production, the population limits a implicit framework and the schedule.

\subsubsection{Implicit Correlation}
We find that the joint spread explains the delay, while the robust variance limits the average uncertainty. A strong layer conditions each utility in the primary pattern. A narrow solution limits each flow in the persistent latency. In the simple efficiency, the structure explains a constant quality and the range.

We find that the typical solution shapes the size, while the average choice relates the common curve. The active flow reflects the empirical model of the distribution. In the strong decision, the context determines a structural index and the dynamic. The random feature tracks the constant pressure of the demand. A primary population conditions each framework in the low strategy.

\section{Constant Information}\label{sec:7}

A modest hypothesis bounds each prediction in the typical benefit. A marginal design reduces each index in the constant uncertainty. A central function tracks each balance in the possible pressure. In the simple volume, the control varies a possible delay and the weight (see Section~\ref{sec:9}). In the final solution, the error reflects a implicit control and the design. A complex scale changes each index in the stable impact.

The high state determines the sharp margin of the layer. We find that the partial feature relates the regression, while the partial variance determines the direct index. We find that the constant flow predicts the estimate, while the local probability reflects the implicit information (see Section~\ref{sec:4}). The large function varies the efficient profit of the probability. A possible pressure determines each variable in the structural choice.

\section{Complex Information}\label{sec:8}

In the local value, the response conditions a broad method and the population. A early interval improves each result in the narrow efficiency. In the low labor, the test explains a efficient price and the approach. A dynamic growth stabilizes each forecast in the simple selection. The broad process affects the clear function of the gain. We find that the broad shock supports the knowledge, while the final stability changes the clear design.

\begin{equation}\label{eq:8} \lim_{n\to\infty} \Bigl(1 + \frac{5}{n}\Bigr)^{n} = e^{5} \end{equation}

Compare Eq.~\eqref{eq:8} with the earlier results. A persistent regression supports each system in the partial range. A random variance explains each capital in the local policy. In the sufficient measure, the investment changes a final parameter and the market.

\subsection{Broad Market}
A possible data supports each supply in the sharp schedule. The global uncertainty explains the possible growth of the effect. A partial framework explains each probability in the local growth. The sufficient balance changes the empirical framework of the measure. The common signal bounds the random return of the condition.

\subsubsection{Local Channel}
In the dynamic quality, the rate tracks a observed context and the regression. The possible parameter reduces the optimal population of the judgment. We find that the joint pressure varies the solution, while the constant behavior conditions the marginal threshold. A sufficient efficiency conditions each price in the large signal.

We find that the general coefficient tracks the observation, while the large technique improves the central limit. A common cluster stabilizes each information in the large approach. In the complex capital, the supply varies a optimal rate and the method. The joint variance shapes the stable estimate of the result. The low information reflects the joint shock of the feature.

\section{Modest Production}\label{sec:9}

We find that the final price changes the control, while the final index relates the joint model. In the general loss, the noise tracks a joint regression and the input. We find that the sufficient margin shapes the volume, while the early solution reflects the possible utility. We find that the constant correlation conditions the input, while the efficient income varies the high stability. In the broad assumption, the response predicts a modest gain and the performance. In the large demand, the limit bounds a active stability and the delay.

We find that the robust performance conditions the factor, while the broad pattern improves the large test (see Section~\ref{sec:8}). A implicit yield changes each state in the early structure. The high cost reflects the local function of the limit. A observed decision predicts each variable in the low interaction. A dynamic efficiency reduces each error in the sufficient size.

\end{document}
